\documentclass[10pt,a4paper]{amsart}
\usepackage[margin=19mm]{geometry}
\usepackage{amsmath,amssymb,mathtools}
\usepackage[scr=rsfs,cal=euler]{mathalfa}
\usepackage{xfrac}
\usepackage[unicode,hidelinks,bookmarks=false]{hyperref}
\newcommand{\ie}{i.e.\ }
\newcommand{\cf}{cf.\ }
\newcommand{\resp}{respectively\ }
\newcommand{\Subsection}{\subsection}
\providecommand{\nobreakdash}{\nobreak\hbox{-}}
\setlength{\emergencystretch}{2em}
\multlinegap0pt
\usepackage{bm}
\usepackage{bbm}
\usepackage{dsfont}
\usepackage{xspace}
\usepackage{cancel}
\usepackage{mathtools}
\usepackage{amsthm}
\makeatletter
\@ifundefined{theorem}{\newtheorem{theorem}{Theorem}}{}
\@ifundefined{lemma}{\newtheorem{lemma}[theorem]{Lemma}}{}
\makeatother
\title[A capacitary approach to Lyapunov-type inequalities for elli]{A capacitary approach to Lyapunov-type inequalities for elliptic problems on weighted graphs}
\author{Mohamed Jleli, Bessem Samet}
\date{Source: arXiv:2607.08405v1 (CC BY 4.0)}
\begin{document}
\maketitle

\noindent\emph{Source and licence.} A capacitary approach to Lyapunov-type inequalities for elliptic problems on weighted graphs. Version arXiv:2607.08405v1 (CC BY 4.0), distributed under the Creative Commons Attribution 4.0 International licence, as recorded in the arXiv licence metadata for this exact version and shown on the abstract page https://arxiv.org/abs/2607.08405v1. Full rights notice: licenses/arxiv-2607.08405-CC-BY-4.0.txt.\par\smallskip

\noindent\textit{Editorial scope.} Counted unit: the Lemma whose source label is \texttt{lem-local-poincare} in the article (source0.tex lines 1618--1630, source label \texttt{lem-local-poincare}), delivered together with the article's own complete original proof of it (source0.tex lines 1632--1792, one \texttt{proof} environment of 161 lines). No mathematics is written, repaired or re-derived: statement and proof are the article's own text line for line. Required context retained uncounted: none. Every definition, display and label of those lines is retained unchanged. Omitted from this package and not redistributed: 1--1617, 1631--1631, 1793--3249. Nothing inside a retained excerpt is omitted or altered: every retained body is delivered exactly as the article's lines are written, so no omission receipt of any kind accompanies this package. Citations: the retained excerpts carry no citation command at all, so no bibliography entry of the article is transcribed and no reference is redistributed. Reference adaptation: none was needed and none was made; every reference inside the delivered unit resolves inside it, so no unresolved reference or citation is produced. Printed slips kept literal: no printed word, symbol, hypothesis or number of the counted statement or of its proof was altered. Layout and class: the article's own class is \texttt{article} with packages including geometry, amsmath, amssymb, amsthm, bm, hyperref, enumitem, xcolor, setspace, float, graphicx, tikz, cite; the wrapper is the lane's pinned \texttt{amsart} editorial wrapper at 10pt on A4, which restates the theorem-like environments the retained text needs and adds the article's own macro definitions that the retained text needs (none), with extra wrapper packages (bm, bbm, dsfont, xspace, cancel, mathtools). The delivered numbering is the unit's own. Assets: no retained span contains a figure environment, a graphics-inclusion command, a PDF-page inclusion, a table or a TikZ picture; no image file is copied into this package, the package declares no asset dependency and no image file is redistributed with it. Licence: version arXiv:2607.08405v1 of this article is distributed under the Creative Commons Attribution 4.0 International licence, as recorded in the arXiv licence metadata for this exact version and shown on the abstract page https://arxiv.org/abs/2607.08405v1. A full rights notice is delivered at licenses/arxiv-2607.08405-CC-BY-4.0.txt. Characters: every retained excerpt is pure ASCII, so the delivered engine, XeLaTeX with its default Latin Modern text font, reports no missing character.
\begin{lemma}[Local Poincar\'e inequality]\label{lem-local-poincare}
Assume that \((V,\omega,\mu)\) satisfies {\rm(G1)} and {\rm(G3)}. Then, for
every \(1<p<\infty\), there exists a constant \(C_P=C_P(p,C_0,C_2)>0\)
such that, for every \(x_0\in V\), every \(R\ge1\), and every
\(u:V\to\mathbb R\),
\begin{equation}\label{local-poincare}
\sum_{x\in B(x_0,R)}\mu(x)|u(x)-u_{B(x_0,R)}|^p
\le
C_P R^p
\sum_{\substack{x,y\in B(x_0,R)\\ y\sim x}}
\omega_{xy}|u(y)-u(x)|^p.
\end{equation}
\end{lemma}

\begin{proof}

Throughout the proof, \(C>0\) denotes a generic constant depending only on \(p,C_0,C_2\).


Let \(x_0\in V\), \(R\ge1\), and \(u:V\to\mathbb R\). Set
\[
B_R=B(x_0,R).
\]
Since the graph is locally finite, \(B_R\) is finite. We may therefore choose
a median \(m\) of \(u\) in \(B_R\), that is,
\[
\mu(\{x\in B_R:\ u(x)>m\})\le \frac{\mu(B_R)}{2},
\qquad
\mu(\{x\in B_R:\ u(x)<m\})\le \frac{\mu(B_R)}{2}.
\]

First, we  prove a weighted \(\ell^p\)-estimate for the positive part \((u-m)_+\). Set
\[
f=(u-m)_+,
\qquad
A_t=\{x\in B_R:\ f(x)>t\},\quad t>0.
\]
Since
\[
A_t\subset \{x\in B_R:\ u(x)>m\},
\]
the median property gives \(\mu(A_t)\le \mu(B_R)/2\). Hence
\[
\min\{\mu(A_t),\mu(B_R\setminus A_t)\}=\mu(A_t).
\]
By {\rm(G3)}, we obtain
\[
\mu(A_t)\le C_2R\,\omega(\partial_{B_R}A_t).
\]
Using the layer-cake formula and the definition of
\(\omega(\partial_{B_R}A_t)\), we get
\[
\begin{aligned}
\sum_{x\in B_R}\mu(x)f(x)^p
&=
p\int_0^\infty t^{p-1}\mu(A_t)\,dt \\
&\le
C_2R
\int_0^\infty p t^{p-1}\omega(\partial_{B_R}A_t)\,dt  \\
&=
C_2R
\sum_{\substack{x,y\in B_R\\ y\sim x}}
\omega_{xy}
\int_0^\infty p t^{p-1}
\mathbf 1_{\{f(y)\le t<f(x)\}}\,dt \\
&\le
C_2R
\sum_{\substack{x,y\in B_R\\ y\sim x}}
\omega_{xy}|f(x)^p-f(y)^p|.
\end{aligned}
\]
Using
\[
|a^p-b^p|
\le
p(a^{p-1}+b^{p-1})|a-b|,
\qquad a,b\ge0,
\]
and H\"older's inequality, we obtain
\[
\sum_{x\in B_R}\mu(x)f(x)^p
\le
C R
\left(
\sum_{\substack{x,y\in B_R\\ y\sim x}}
\omega_{xy}|f(y)-f(x)|^p
\right)^{1/p} 
\left(
\sum_{\substack{x,y\in B_R\\ y\sim x}}
\omega_{xy}
\bigl(f(x)^{p-1}+f(y)^{p-1}\bigr)^{\frac{p}{p-1}}
\right)^{\frac{p-1}{p}}.
\]
Moreover,
\[
\bigl(f(x)^{p-1}+f(y)^{p-1}\bigr)^{\frac{p}{p-1}}
\le
C\bigl(f(x)^p+f(y)^p\bigr).
\]
Using {\rm(G1)}, namely the boundedness of the edge weights and degrees
together with the lower bound on \(\mu\), we obtain
\[
\sum_{\substack{x,y\in B_R\\ y\sim x}}
\omega_{xy}\bigl(f(x)^p+f(y)^p\bigr)
\le
C
\sum_{x\in B_R}\mu(x)f(x)^p.
\]
Therefore,
\[
\sum_{x\in B_R}\mu(x)f(x)^p
\le
C R
\left(
\sum_{\substack{x,y\in B_R\\ y\sim x}}
\omega_{xy}|f(y)-f(x)|^p
\right)^{1/p}
\left(
\sum_{x\in B_R}\mu(x)f(x)^p
\right)^{\frac{p-1}{p}}.
\]
Since \(s\mapsto (s-m)_+\) is \(1\)-Lipschitz, we have
\[
|f(y)-f(x)|\le |u(y)-u(x)|,
\qquad x,y\in B_R.
\]
Consequently,
\begin{equation}\label{chang1}
\sum_{x\in B_R}\mu(x)f(x)^p
\le
C R^p
\sum_{\substack{x,y\in B_R\\ y\sim x}}
\omega_{xy}|u(y)-u(x)|^p.
\end{equation}

Applying the same argument to \(g=(m-u)_+\), we obtain
\begin{equation}\label{chang2}
\sum_{x\in B_R}\mu(x)g(x)^p
\le
C R^p
\sum_{\substack{x,y\in B_R\\ y\sim x}}
\omega_{xy}|u(y)-u(x)|^p.
\end{equation}
Since \(|u-m|^p=f^p+g^p\), it follows from \eqref{chang1} and \eqref{chang2} that
\begin{equation}\label{chang3}
\sum_{x\in B_R}\mu(x)|u(x)-m|^p
\le
C R^p
\sum_{\substack{x,y\in B_R\\ y\sim x}}
\omega_{xy}|u(y)-u(x)|^p.
\end{equation}

It remains to pass from the median to the average. By Jensen's inequality,
\[
|u_{B_R}-m|^p
\le
\frac{1}{\mu(B_R)}
\sum_{x\in B_R}\mu(x)|u(x)-m|^p.
\]
Hence
\begin{equation}\label{chang4}
\begin{aligned}
\sum_{x\in B_R}\mu(x)|u(x)-u_{B_R}|^p
&\le
2^{p-1}\sum_{x\in B_R}\mu(x)|u(x)-m|^p
+
2^{p-1}\mu(B_R)|u_{B_R}-m|^p  \\
&\le
2^p
\sum_{x\in B_R}\mu(x)|u(x)-m|^p.
\end{aligned}
\end{equation}
Combining \eqref{chang3} and \eqref{chang4}, we
obtain \eqref{local-poincare} for some \(C_P=C_P(p,C_0,C_2)>0\).
\end{proof}

\end{document}
